\subsection{平展映射的连续截面的构造}

定理 1. --- 设 X、Y、Z 为拓扑空间，\(f: \mathrm{Z} \to \mathrm{X} \times \mathrm{Y}\) 为一个分离平展态射，并设 \(y_0\) 为 Y 的一点。设 \(s: X \times Y \to Z\) 为 \(f\) 的一个截面。假定 \(s\) 在 \(X \times \{y_0\}\) 上的限制是连续的，并且 \(s\) 在 \(\{x\} \times Y\) （对每个 \(x \in X\) ）上的限制也都是连续的。若空间 \(Y\) 是连通且局部连通的 (TG, I, p. 84, def. 4)，则映射 \(s\) 是连续的。

引理. --- 设 U 为 X 的一个开子集，V 为 Y 的一个连通开子集，并设 \((x_{1},y_{1})\in\mathrm{U}\times\mathrm{V}\) 。假定 s 在 \(U\times\{y_{1}\}\) 上的限制是连续的。设 \(\sigma\) 为 f 在 \(U\times V\) 上的一个连续截面，使得 \(\sigma(x_{1},y_{1})=s(x_{1},y_{1})\) 。则存在一个邻域 \(U'\) （\(x_{1}\) 的），使得 \(s=\sigma\) 在 \(U'\times V\) 上成立。特别地，映射 s 在 \((x_{1},y_{1})\) 的一个邻域中是连续的。

由于 \(s\) 在 \(\mathrm{U} \times \{y_1\}\) 上的限制是连续的，且映射 \(f\) 是平展的，故存在一个邻域 \(\mathrm{U}'\) （\(x_1\) 的），使得 \(s(x, y_1) = \sigma(x, y_1)\) 对一切 \(x \in \mathrm{U}'\) 成立 (I, p. 34, prop. 11, b))。设 \(x \in \mathrm{U}'\) ；由于 \(s\) 在 \(\{x\} \times \mathrm{V}\) 上的限制是连续的，且映射 \(f\) 是平展且分离的，由 I, p. 34 的推论 1 得出，\(\sigma(x, y) = s(x, y)\) 对一切 \(y \in \mathrm{V}\) 成立。于是，\(s\) 与 \(\sigma\) 在 \(\mathrm{U}' \times \mathrm{V}\) 上重合。

现在来证明定理。首先证明映射 s 在 \(X \times \{y_{0}\}\) 的每一点的一个邻域中是连续的。设 \(x_{0} \in X\) ；可以选取 X 中 \(x_{0}\) 的一个开邻域 U、Y 中 \(y_{0}\) 的一个连通开邻域 V，以及 f 的一个连续截面 \(\sigma: U \times V \to Z\) ，使得 \(\sigma(x_{0}, y_{0}) = s(x_{0}, y_{0})\) 。由上述引理，s 在 \((x_{0}, y_{0})\) 的一个邻域中是连续的。

我们现证明映射 \(s\) 在 \(\mathrm{X} \times \mathrm{Y}\) 的每一点处都是连续的。对 \(x_0 \in \mathrm{X}\) ，用 \(\mathrm{C}(x_0)\) 表示点 \(y \in \mathrm{Y}\) （使 \(s\) 在 \((x_0, y)\) 的一个邻域中连续者）所成的集。按定义，集 \(\mathrm{C}(x_0)\) 在 Y 中是开的，且由前所证它含有 \(y_0\) 。

我们来证明它在 Y 中是闭的。考虑 Y 中的一点 \(y_{1}\) （属于 \(\mathrm{C}(x_{0})\) 的闭包者）、X 中 \(x_{0}\) 的一个开邻域 U、Y 中 \(y_{1}\) 的一个开邻域 V，以及 f 的一个连续截面 \(\sigma\colon U\times V\to Z\) ，使得 \(\sigma(x_{0},y_{1})=s(x_{0},y_{1})\) 。由于 s 在 \(\{x_{0}\}\times V\) 上的限制连续且 f 是平展的，故存在一个开邻域 \(V^{\prime}\) （\(y_{1}\) 的，在 Y 中），使得 \(\sigma(x_{0},y)=s(x_{0},y)\) 对一切 \(y\in V^{\prime}\) 成立 (I, p. 34 的命题 11)。由于 Y 是局部连通的，可以假设 \(V^{\prime}\) 是连通的。由于 \(y_{1}\) 属于 \(\mathrm{C}(x_{0})\) 的闭包，集 \(\mathrm{C}(x_{0})\cap\mathrm{V}^{\prime}\) 非空；设 \(y_{2}\) 为 \(\mathrm{C}(x_{0})\cap\mathrm{V}^{\prime}\) 的一点。把引理应用于 \((x_{0},y_{2})\) ，便存在一个邻域 \(U^{\prime}\) （\(x_{0}\) 的，含于 U 中），使得 \(s=\sigma\) 在 \(U^{\prime} \times V^{\prime}\) 上成立。这就证明了 \(\mathrm{C}(x_{0})\) 是闭的，因为 \(U^{\prime} \times V^{\prime}\) 是 \((x_{0}, y_{1})\) 的一个邻域。由于 Y 是连通的，我们有 \(\mathrm{C}(x_{0}) = \mathrm{Y}\) ，且这对每个 \(x_{0} \in X\) 都成立，由此即得定理。

推论 1. --- 设 X、Y、E、B 为拓扑空间。设 \(p\colon E \to B\) 为一个平展且分离的映射，\(h\colon X \times Y \to B\) 为一个连续映射，并设 \(y_0\) 为 Y 的一点。设 \(g\colon X \times Y \to E\) 为满足 \(p \circ g = h\) 的一个映射。假设 \(g\) 在 \(X \times \{y_0\}\) 上的限制以及，对每个 \(x \in X\) ，在 \(\{x\} \times Y\) 上的限制都是连续的。若空间 Y 是连通且局部连通的，则映射 \(g\) 是连续的。

设 Z 为纤维积 \((\mathrm{X} \times \mathrm{Y}) \times_{\mathrm{B}} \mathrm{E}\) ，并设 \(f: Z \to X \times Y\) 与 \(h': Z \to E\) 为典范投影。映射 f 是平展的 (I, p. 31, prop. 8) 且是分离的 (I, p. 27, prop. 4)。映射 \(s: X \times Y \to Z\) （由 \(s(x, y) = ((x, y), g(x, y))\) 定义）是 f 的一个截面，它满足定理 1 的诸假设。因此它是连续的，\(g = h' \circ s\) 亦然。

注. --- 若在定理 1 中不假设空间 Y 是局部连通的，则定理的结论不再必然成立 (I, p. 141, exerc. 7)。

定理 2. --- 设 B 为一个拓扑空间，A 为 B 的一个子空间。假设下列假设之一成立：

\begin{enumerate}
\def\labelenumi{(\roman{enumi})}
\item
  A 具有由仿紧邻域组成的基本邻域系；
\item
  B 是仿紧的且 A 是闭的；
\item
  B 是可度量化的；
\item
  A 是紧的，且 A 中任意两个不同的点在 B 中有不相交的邻域。
\end{enumerate}

那么，平展映射 \(p\colon E \to B\) 在 A 上的每个连续截面都可以扩张为 \(p\) 在 A 的一个邻域上的连续截面。

考虑数对 (B, A) 的下列性质 (PCV)：

(PCV) 对由 B 的开子集组成的 A 的每个覆盖 \((\mathrm{U}_{i})_{i\in\mathrm{I}}\) ，存在 A 的一个邻域 V 和覆盖 V 的、由 V 的闭子集组成的局部有限族 \((\mathrm{F}_{j})_{j\in\mathrm{J}}\) ，使得每个 \(F_{j}\) 都含于某个 \(U_{i}\) 之中。

定理 2 由下文的引理 1 与引理 3 得出。

引理 1. --- 定理 2 的性质 (i) 至 (iv) 中的每一条都蕴涵性质 (PCV)。

设 \((\mathrm{U}_{i})_{i\in\mathrm{I}}\) 为由 B 的开子集组成的 A 的一个覆盖。在假设 (i) 之下，存在 A 的一个仿紧邻域 V，它含于诸 \(U_{i}\) 的并，存在空间 V 的一个开覆盖 \((\mathrm{U}_{j}^{\prime})_{j\in\mathrm{J}}\) ，它是局部有限的且加细于覆盖 \((\mathrm{V}\cap\mathrm{U}_{i})_{i\in\mathrm{I}}\) ，还存在空间 V 的一个开覆盖 \((\mathrm{U}_{j}^{\prime\prime})_{j\in\mathrm{J}}\) ，使得对每个 \(j\in J\) ，闭包 \(F_{j}\) （\(U_{j}^{\prime\prime}\) 在 V 中的）含于 \(U_{j}^{\prime}\) (TG, IX, p. 49, prop. 4 及 p. 48, th. 3 的 cor. 1)。故在此情形下 (PCV) 成立。

由 TG, IX, p. 50 的推论 2，条件 (ii) 蕴涵条件 (i)。同样，条件 (iii) 蕴涵条件 (i)：事实上，若 B 是可度量化的，则 A 的每个邻域都是可度量化的，从而是仿紧的 (TG, IX, p. 51, th. 4)。

最后假设条件 (iv) 成立，我们来证明性质 (PCV) 成立。由于 A 是紧的，只需考虑有限覆盖 \((\mathrm{U}_{i})_{0\leqslant i\leqslant n}\) 的情形。于是用归纳法构造 B 的开集 \(V_{0},\ldots,V_{n}\) ，使其对 \(0\leqslant i\leqslant n\) 满足下列条件：

\(\alpha)\mathrm{A}\subset \mathrm{V}_{0}\cup \dots \cup \mathrm{V}_{i}\cup \mathrm{U}_{i + 1}\cup \dots \cup \mathrm{U}_{n};\)

\(\beta)\overline{\mathrm{V}_{i}}\cap \mathrm{A}\subset \mathrm{U}_{i}.\)

设 \(V_{0},\ldots,V_{r-1}\) 已经构造出来，且对 \(0\leqslant i\leqslant r-1\) 满足上述条件，而来构造 \(V_{r}\) 。集合 \(K=A\cap C U_{r}\) 与 \(L=A\cap C(V_{0}\cup\cdots\cup V_{r-1}\cup U_{r+1}\cup\cdots\cup U_{n})\) 在 A 中是闭的，因而是紧的。根据 \((\alpha)\) ，它们不相交。由假设，对 L 的每一点 a 与 K 的每一点 b，存在 a 与 b 在 B 中的不相交邻域。由下文的引理 2，存在 B 中的开集 \(V_{r}\) 与 W，二者不相交，且 \(L\subset V_{r}\) 与 \(K\subset W\) 。由包含关系 \(L\subset V_{r}\) 与 \(A\subset V_{0}\cup\cdots\cup V_{r-1}\cup U_{r}\cup\cdots\cup U_{n}\) 以及 L 的定义，可对 i=r 推出 \((\alpha)\) 。另一方面，有 \(\overline{V_{r}}\cap K=\varnothing\) ，从而 \(\overline{V_{r}}\cap A\subset U_{r}\) 。

集合 \(M = \bigcup_{0 \leqslant i \leqslant n} (\overline{\mathrm{V}_{i}} \cap \mathbb{C}\mathrm{U}_{i})\) 是闭的，且由 \(\beta\) 知它与 A 不相交。由 \(\alpha\) ，集合 \(V = (\bigcup_{0 \leqslant i \leqslant n} \overline{\mathrm{V}_{i}}) - M\) 是 A 在 B 中的一个邻域。对 \(i = 0, \ldots, n\) ，令 \(F_{i} = V \cap \overline{V_{i}}\) ；它是 V 的一个闭子集，含于 \(U_{i}\) 。族 \((F_{i})_{0 \leqslant i \leqslant n}\) 是 V 的一个覆盖，故性质 (PCV) 成立。

引理 2. --- 设 B 为一个拓扑空间，并设 K 与 L 为 B 的拟紧子集。假设对 K 的每一点 a 与 L 的每一点 b，存在 a 与 b 在 B 中的不相交邻域。那么存在 B 中两个不相交的开集 U 与 V，使得 K ⊂ U 且 L ⊂ V。

设 \(a\) 为 K 的一点。对每个 \(b \in L\) ，设 \(U_{a,b}\) 与 \(V_{a,b}\) 分别为 \(a\) 与 \(b\) 的不相交开邻域。固定 \(a\) ，族 \((V_{a,b})_{b \in L}\) 是 L 的一个开覆盖。由于 L 是拟紧的，存在一个有限族 \(T_a \subset L\) ，使得 \(V_b\) （诸 \(V_{a,b}\) 对 \(b \in T_a\) 者的并）包含 L。交 \(U_a\) （诸 \(U_{a,b}\) 对 \(b \in T_a\) 者之交）是 \(a\) 的一个开邻域，因为 \(T_a\) 是有限的；此外，\(U_a\) 与 \(V_a\) 不相交。诸 \((U_a)_{a \in A}\) 构成 K 的一个开覆盖。由于 K 是拟紧的，存在一个有限族 \(S \subset K\) ，使得 \(U = \bigcup_{a \in S} U_a\) 包含 K。于是，\(V = \bigcap_{a \in S} V_a\) 是 B 的一个包含 L 的开子集。引理得证。

引理 3. --- 设 B 为一个拓扑空间，A 为 B 的一个子空间，使得数对 (B,A) 满足 I, p. 37 的性质 (PCV)。那么，平展映射 \(p\colon E \to B\) 在 A 上的每个连续截面都可以扩张为 \(p\) 在 A 的一个邻域上的连续截面。

设 \(p \colon \mathrm{E} \to \mathrm{B}\) 为一个平展映射，\(s \colon \mathrm{A} \to \mathrm{E}\) 为 \(p\) 在 A 上的一个连续截面。对 A 的每一点 \(a\) ，存在一个开邻域 \(\mathrm{U}_a\) （\(a\) 的）和一个连续截面 \(s_a \colon \mathrm{U}_a \to \mathrm{E}\) ，它与 \(s\) 在 \(\mathrm{U}_a \cap \mathrm{A}\) 上重合 (I, p. 34, prop. 10 及 I, p. 34, prop. 11)。族 \((\mathrm{U}_a)_{a \in \mathrm{A}}\) 是由 B 的开集组成的 A 的一个覆盖。设 V 为 A 在 B 中的一个邻域，\((\mathrm{F}_j)_{j \in \mathrm{J}}\) 为覆盖 V 的、由 V 的闭子集组成的局部有限族，并设 \(a \colon \mathrm{J} \to \mathrm{A}\) 为一个映射，使得 \(\mathrm{F}_j\) 含于 \(\mathrm{U}_{a(j)}\) （对一切 \(j \in \mathrm{J}\) ；性质 (PCV)）。用 \(s_j\) 表示 \(s_{a(j)}\) 在 \(\mathrm{F}_j\) 上的限制。设 W 为满足下列条件的点 \(b \in \mathrm{V}\) 所成的集：\(s_j(b) = s_k(b)\) 对每对 \((j, k) \in \mathrm{J} \times \mathrm{J}\) （使 \(b \in \mathrm{F}_j \cap \mathrm{F}_k\) 者）都成立。我们定义截面 \(s' \colon \mathrm{W} \to \mathrm{E}\) 如下：\(s'(b) = s_j(b)\) （若 \(b \in \mathrm{F}_j\) ）。我们有 \(\mathrm{A} \subset \mathrm{W}\) 且 \(s' | \mathrm{A} = s\)。由 I, p. 19 的 prop. 4，截面 \(s'\) 是连续的。

尚须证明 W 是 A 在 B 中的一个邻域。对每对 \((j,k)\in\mathrm{J}\times\mathrm{J}\) ，用 \(T_{jk}\) 表示点 \(b\in F_{j}\cap F_{k}\) （满足 \(s_{j}(b)\neq s_{k}(b)\) 者）所成的集。集 \(T_{jk}\) 在 \(F_{j}\cap F_{k}\) 中是闭的 (I, p. 34 的 prop. 11, b))，从而在 V 中是闭的，且诸集 \(T_{jk}\) 构成 V 的子集的一个局部有限族。于是族 \((T_{jk})\) 之并 T 是 V 中的一个闭集 (TG, I, p. 6, prop. 4)。而按定义，W 是 T 在 V 中的余集。因此它是 A 在 V 中的一个邻域，从而也是 A 在 B 中的邻域。

